% ====================================================================
% fig_ff1_5_kinetics.tex — [FF1 후보 5] 핵생성 장벽과 spinodal 성장 밴드
%   대상 문서: appendix_phase_separation (라이브러리: arrows.meta·calc·backgrounds)
%   제안 배치: §A.6 두 소절 뒤(두 경로 대조표 앞)
%   좌표 생성: eval_ff1_coords.py §P5 — (a) 식 eq:app-cnt 를 r*=2γ/|Δg_v|·
%   ΔG*=16πγ³/3Δg_v² 로 무차원화한 ΔG/ΔG*=3s²−2s³ (s=r/r*) 과 그 두 항,
%   (b) 식 eq:app-ch-R 을 k_c=√(−f''/2κ)·M|f''|k_c² 로 무차원화한 성장률.
% ====================================================================
\begin{figure}[t]
\centering
\begin{tikzpicture}[>=Latex]
% =============== (a) nucleation barrier ===============
\begin{scope}[x=3.30cm,y=0.56cm]
\draw[->] (0,0) -- (2.02,0) node[below,font=\scriptsize] {$r/r^\ast$};
\draw[->] (0,-2.9) -- (0,3.4) node[above,font=\scriptsize] {$\Delta G(r)/\Delta G^\ast$};
\foreach \xx in {0.5,1.0,1.5} {\draw (\xx,0.06) -- (\xx,-0.06) node[below,font=\scriptsize] {\xx};}
\foreach \yy in {-2,-1,1,2,3} {\draw (0.012,\yy) -- (-0.012,\yy) node[left,font=\scriptsize] {\yy};}
% surface term 3s^2 (dotted)
\draw[densely dotted,thick] plot[smooth] coordinates {(0,0) (0.1,0.03) (0.2,0.12) (0.3,0.27)
  (0.4,0.48) (0.5,0.75) (0.6,1.08) (0.7,1.47) (0.8,1.92) (0.9,2.43) (1,3)};
\node[font=\scriptsize,anchor=west] at (0.62,2.30) {계면 $4\pi r^2\gamma$};
\draw[->,thin,black!60] (0.79,2.20) -- (0.87,2.12);
% volume term -2s^3 (dashed)
\draw[densely dashed,thick] plot[smooth] coordinates {(0,0) (0.1,-0.002) (0.2,-0.016) (0.3,-0.054)
  (0.4,-0.128) (0.5,-0.25) (0.6,-0.432) (0.7,-0.686) (0.8,-1.024) (0.9,-1.458) (1,-2) (1.1,-2.662)};
\node[font=\scriptsize,anchor=east,align=right] at (0.98,-1.35) {체적 $\tfrac43\pi r^3\Delta g_v$\\($\Delta g_v<0$)};
% total (thick)
\draw[very thick] plot[smooth] coordinates {(0,0) (0.1,0.028) (0.2,0.104) (0.3,0.216) (0.4,0.352)
  (0.5,0.5) (0.6,0.648) (0.7,0.784) (0.8,0.896) (0.9,0.972) (1,1) (1.1,0.968) (1.2,0.864)
  (1.3,0.676) (1.4,0.392) (1.5,0) (1.6,-0.512) (1.7,-1.156) (1.8,-1.944) (1.9,-2.888)};
\node[font=\scriptsize,anchor=west] at (1.28,1.15) {합 $\Delta G(r)$};
\draw[->,thin,black!60] (1.30,1.00) -- (1.24,0.82);
% critical nucleus markers
\draw[densely dotted,black!60] (1,0) -- (1,1);
\draw[densely dotted,black!60] (0,1) -- (1,1);
\fill (1,1) circle (1.1pt);
\node[font=\scriptsize,anchor=south west] at (0.02,1.02) {$\Delta G^\ast$};
\node[font=\scriptsize,anchor=north] at (1,-0.30) {$r^\ast$};
% fate arrows
\draw[->,semithick] (0.62,-0.68) -- (0.30,-0.68);
\node[font=\tiny,anchor=west] at (0.10,-1.05) {$r<r^\ast$: 소멸};
\draw[->,semithick] (1.38,1.62) -- (1.72,1.62);
\node[font=\tiny,anchor=west] at (1.30,2.00) {$r>r^\ast$: 성장};
\node[font=\scriptsize] at (1.0,-3.55) {(a) 준안정: 핵생성 장벽 (식~\eqref{eq:app-cnt}$\cdot$\eqref{eq:app-rstar})};
\end{scope}
% =============== (b) spinodal decomposition growth band ===============
\begin{scope}[xshift=8.35cm,x=4.75cm,y=2.75cm]
% growth band shading (under growth curve, 0<k<k_c)
\fill[black!10] plot[smooth] coordinates {(0,0) (0.1,0.0099) (0.2,0.0384) (0.3,0.0819)
  (0.4,0.1344) (0.5,0.1875) (0.6,0.2304) (0.7071,0.25) (0.8,0.2304) (0.9,0.1539) (1,0)} -- cycle;
\draw[->] (0,0) -- (1.42,0) node[below,font=\scriptsize] {$k/k_c$};
\draw[->] (0,-1.12) -- (0,0.42) node[above,font=\scriptsize] {$R(k)\big/\big[M|f''|k_c^2\big]$};
\foreach \xx in {0.5,1.0} {\draw (\xx,0.012) -- (\xx,-0.012) node[below,font=\scriptsize] {\xx};}
\foreach \yy/\lab in {0.25/{$\tfrac14$},-0.5/{$-\tfrac12$},-1/{$-1$}}
  {\draw (0.006,\yy) -- (-0.006,\yy) node[left,font=\scriptsize] {\lab};}
% f''<0 growth curve
\draw[very thick] plot[smooth] coordinates {(0,0) (0.1,0.0099) (0.2,0.0384) (0.3,0.0819)
  (0.4,0.1344) (0.5,0.1875) (0.6,0.2304) (0.7071,0.25) (0.8,0.2304) (0.9,0.1539) (1,0)
  (1.1,-0.2541) (1.2,-0.6336) (1.28,-1.02)};
% f''>0 all-mode decay (mirror normalization with +|f''|)
\draw[dashed,black!55] plot[smooth] coordinates {(0,0) (0.1,-0.0101) (0.2,-0.0416) (0.3,-0.0981)
  (0.4,-0.1856) (0.5,-0.3125) (0.6,-0.4896) (0.7071,-0.75) (0.8,-1.0496)};
\node[font=\scriptsize,black!55,anchor=west,align=left] at (0.72,-0.62)
  {$f''>0$ (spinodal 밖):\\ 전-모드 감쇠 $R<0$};
% fastest mode marker
\draw[densely dotted,black!60] (0.7071,0) -- (0.7071,0.25);
\fill (0.7071,0.25) circle (1.1pt);
\node[font=\scriptsize,anchor=south,align=center] at (0.62,0.29)
  {$k_m{=}k_c/\sqrt2$: $R_{\max}{=}\tfrac14$};
% band annotation
\node[font=\scriptsize,anchor=south west] at (0.035,0.025) {성장 밴드 $0<k<k_c$};
\node[font=\scriptsize,anchor=north] at (1.0,-0.05) {$k_c$};
\node[font=\scriptsize] at (0.71,-1.32) {(b) 불안정: spinodal 분해 성장률 (식~\eqref{eq:app-ch-R})};
\end{scope}
\end{tikzpicture}
\caption{준안정$\cdot$불안정 영역의 두 분해 동역학 --- 좌표는 두 식을 무차원화해 그대로 수치 평가한
값이다. (a) 구형 핵의 자유에너지 식~\eqref{eq:app-cnt} 를 임계 핵 $r^\ast=2\gamma/|\Delta g_v|$$\cdot$
장벽 $\Delta G^\ast=16\pi\gamma^3/3\Delta g_v^2$(식~\eqref{eq:app-rstar})로 접으면 보편형
$\Delta G/\Delta G^\ast=3s^2-2s^3$($s{=}r/r^\ast$): 계면 항(점선)이 앞서고 체적 항(파선)이 뒤에 이겨
$r^\ast$ 에서 장벽 $\Delta G^\ast$ 를 만들며, $r<r^\ast$ 핵은 녹고 $r>r^\ast$ 핵만 자란다 --- 핵생성
속도가 $\exp(-\Delta G^\ast/k_BT)$ 로 지수 억제되는 까닭이다. (b) 선형화 성장률 식~\eqref{eq:app-ch-R}
을 $k_c=\sqrt{-f''/2\kappa}$ 와 $M|f''|k_c^2$ 로 접은 보편형: $f''<0$(spinodal 안)이면 장파장 밴드
$0<k<k_c$(음영) 전체가 장벽 없이 지수 성장하고 $k_m=k_c/\sqrt2$ 의 조성 물결이 가장 빨리 증폭되며
($R_{\max}=M|f''|k_c^2/4=Mf''^2/8\kappa$), $f''>0$ 이면(파선) 모든 모드가 감쇠해 균일상이 유지된다.}
\label{fig:cand-ff1-5}
\end{figure}
